\section*{Bloque B. Aplicaciones cuantitativas y optimización}

\noindent\textit{Preguntas 31 a 65. Requieren plantear y resolver el problema de optimización del consumidor con distintas especificaciones funcionales, así como calcular efectos precio, efectos ingreso, medidas de bienestar y elasticidades.}

\q{Un consumidor tiene $U(x,y)=x^{0.5}y^{0.5}$, ingreso $I=200$, $p_x=4$ y $p_y=2$. La cantidad óptima de $x$ es:}
\begin{opts}
  \item 12.5
  \item 20
  \item 50
  \item 25
\end{opts}

\q{Con $U(x,y)=x^{1/3}y^{2/3}$, $I=300$, $p_x=5$ y $p_y=10$, el gasto óptimo en $y$ es:}
\begin{opts}
  \item 200
  \item 100
  \item 150
  \item 250
\end{opts}

\q{Con $U(x,y)=4x+3y$, $p_x=2$, $p_y=3$ e $I=100$, el consumidor elige:}
\begin{opts}
  \item $x=0$, $y=33.33$
  \item $x=25$, $y=16.67$
  \item $x=50$, $y=0$
  \item $x=20$, $y=20$
\end{opts}

\q{Con $U(x,y)=\min(2x,y)$, $p_x=3$, $p_y=2$ e $I=140$, la utilidad óptima es:}
\begin{opts}
  \item 20
  \item 40
  \item 30
  \item 70
\end{opts}

\q{Con $U(x,y)=\min(x,2y)$, $p_x=4$, $p_y=2$ e $I=120$, la cantidad óptima de $x$ es:}
\begin{opts}
  \item 12
  \item 20
  \item 30
  \item 24
\end{opts}

\q{Para $U(x,y)=(x^{0.5}+y^{0.5})^2$, la elasticidad de sustitución es:}
\begin{opts}
  \item 2
  \item 0.5
  \item 1
  \item infinita
\end{opts}

\q{Para $U(x,y)=(x^{-1}+y^{-1})^{-1}$, la elasticidad de sustitución es:}
\begin{opts}
  \item 0
  \item 1
  \item 0.5
  \item 2
\end{opts}

\q{Con $U(x,y)=\ln x+y$, $p_x=2$, $p_y=1$ e $I=20$, la demanda óptima de $x$ es:}
\begin{opts}
  \item 0.25
  \item 0.5
  \item 1
  \item 2
\end{opts}

\q{Con $U(x,y)=x+\ln y$, $p_x=1$, $p_y=2$ e $I=10$, la demanda óptima de $y$ es:}
\begin{opts}
  \item 1
  \item 2
  \item 5
  \item 0.5
\end{opts}

\q{Si $U(x,y)=x^{0.3}y^{0.6}$, la participación del gasto en $x$ es:}
\begin{opts}
  \item 33.3\%
  \item 30\%
  \item 50\%
  \item 66.7\%
\end{opts}

\q{Para una función Cobb-Douglas $U=x^ay^{1-a}$, la elasticidad precio de la demanda de $x$ es:}
\begin{opts}
  \item 0
  \item 1
  \item $-1$
  \item $-a$
\end{opts}

\q{Para la misma función Cobb-Douglas, la elasticidad cruzada de la demanda de $x$ respecto de $p_y$ es:}
\begin{opts}
  \item $-1$
  \item 0
  \item 1
  \item $a$
\end{opts}

\q{Con preferencias Stone-Geary $U=(x-10)(y-5)$, $p_x=2$, $p_y=4$ e $I=100$, la demanda óptima de $x$ es:}
\begin{opts}
  \item 20
  \item 30
  \item 37.5
  \item 25
\end{opts}

\q{En el problema anterior, el gasto óptimo en $y$ es:}
\begin{opts}
  \item 50
  \item 40
  \item 60
  \item 75
\end{opts}

\q{Con $U=xy$, $I=120$ y $p_x=p_y=1$, el precio de $x$ sube a 2. El efecto sustitución de Slutsky sobre $x$ es:}
\begin{opts}
  \item $-30$
  \item $-10$
  \item $-15$
  \item $0$
\end{opts}

\q{En el mismo caso anterior, el efecto ingreso sobre $x$ es:}
\begin{opts}
  \item $-30$
  \item $-15$
  \item $-10$
  \item $0$
\end{opts}

\q{Con $U=\sqrt{xy}$, $I=100$ y $p_x=p_y=1$, el precio de $x$ sube a 4. La variación compensatoria (CV) es:}
\begin{opts}
  \item 50
  \item 75
  \item 125
  \item 100
\end{opts}

\q{En el mismo caso anterior, la variación equivalente (EV), expresada como disposición a aceptar, es:}
\begin{opts}
  \item 50
  \item 25
  \item 100
  \item 150
\end{opts}

\q{La demanda inversa es $P=100-2Q$. Si el precio sube de 60 a 80, la pérdida de excedente del consumidor es:}
\begin{opts}
  \item 100
  \item 200
  \item 300
  \item 400
\end{opts}

\q{Con $U=2\sqrt{x}+y$, $p_x=p_y=1$ e $I=100$, el precio de $x$ sube a 4. La variación compensatoria es:}
\begin{opts}
  \item 0.5
  \item 0.75
  \item 1
  \item 2
\end{opts}

\q{Para $Q=100-2P$, la elasticidad precio en $P=20$ es:}
\begin{opts}
  \item $-0.33$
  \item $-1.50$
  \item $-2.00$
  \item $-0.67$
\end{opts}

\q{Para $Q=100P^{-2}$, la elasticidad precio es:}
\begin{opts}
  \item $-2$
  \item $-0.5$
  \item $-1$
  \item $0$
\end{opts}

\q{Si la elasticidad precio es $-0.5$ y el precio aumenta 10\%, la cantidad demandada:}
\begin{opts}
  \item cae 10\%
  \item cae 0.5\%
  \item cae 5\%
  \item aumenta 5\%
\end{opts}

\q{Si la elasticidad ingreso es 0.8 y el ingreso aumenta 10\%, la demanda:}
\begin{opts}
  \item aumenta 10\%
  \item aumenta 8\%
  \item cae 8\%
  \item no cambia
\end{opts}

\q{La elasticidad cruzada del agua potable respecto del precio del agua embotellada es 0.3. Si el precio del agua embotellada sube 10\%, la demanda de agua potable:}
\begin{opts}
  \item cae 3\%
  \item no cambia
  \item aumenta 30\%
  \item aumenta 3\%
\end{opts}

\q{Con $U=2x+y$, $p_x=3$, $p_y=2$ e $I=60$, el consumidor elige:}
\begin{opts}
  \item $x=20$, $y=0$
  \item $x=0$, $y=30$
  \item $x=10$, $y=15$
  \item $x=12$, $y=12$
\end{opts}

\q{Con $U=\min(x,3y)$, $p_x=2$, $p_y=1$ e $I=70$, la cantidad óptima de $x$ es:}
\begin{opts}
  \item 10
  \item 20
  \item 30
  \item 35
\end{opts}

\q{Con $U=\sqrt{xy}$, $p_x=2$, $p_y=8$ e $I=160$, la RMS en el óptimo es:}
\begin{opts}
  \item $1/2$
  \item $1/4$
  \item $2$
  \item $4$
\end{opts}

\q{Con $U=x^2y$, $p_x=3$, $p_y=1$ e $I=90$, la demanda óptima de $x$ es:}
\begin{opts}
  \item 10
  \item 15
  \item 30
  \item 20
\end{opts}

\q{La elasticidad de sustitución de una función de utilidad Leontief es:}
\begin{opts}
  \item 0
  \item 0.5
  \item 1
  \item infinita
\end{opts}

\q{Para $U=\sqrt{xy}$, la función de utilidad indirecta es:}
\begin{opts}
  \item $v(p,I)=\dfrac{I}{\sqrt{p_xp_y}}$
  \item $v(p,I)=\dfrac{2I}{\sqrt{p_xp_y}}$
  \item $v(p,I)=\dfrac{I}{2\sqrt{p_xp_y}}$
  \item $v(p,I)=\dfrac{I}{p_x+p_y}$
\end{opts}

\q{Para $U=\sqrt{xy}$, $p_x=4$ y $p_y=1$, el gasto mínimo para alcanzar $u=10$ es:}
\begin{opts}
  \item 20
  \item 40
  \item 80
  \item 100
\end{opts}

\q{Si $x=I/(2p_x)$, un aumento del ingreso de 100 a 200, con $p_x=5$, eleva $x$ en:}
\begin{opts}
  \item 5
  \item 15
  \item 20
  \item 10
\end{opts}

\q{Para $U=xy$, con $p_y$ e $I$ constantes, la curva de precio-consumo es:}
\begin{opts}
  \item horizontal.
  \item vertical.
  \item una recta que pasa por el origen.
  \item una curva con pendiente positiva.
\end{opts}

\q{Para $U=xy$, la curva de ingreso-consumo es:}
\begin{opts}
  \item horizontal.
  \item vertical.
  \item una recta que pasa por el origen.
  \item una curva de pendiente negativa.
\end{opts}
